\documentclass[a4paper,12pt]{article}
\usepackage[UTF8]{ctex}
\usepackage{geometry}
\geometry{left=2.5cm,right=2.5cm,top=2.5cm,bottom=2.5cm}
\usepackage{fancyhdr}
\usepackage{graphicx}
\usepackage{longtable}
\usepackage{booktabs}
\usepackage{xcolor}
\usepackage{listings}
\usepackage{hyperref}
\usepackage{titlesec}
\usepackage{caption}

% 封面命令
\newcommand{\doccover}[1]{
\begin{titlepage}
\centering
\vspace*{3cm}
{\Huge\bfseries #1\par}
\vspace{1.5cm}
{\LARGE 在线论文文档自动生成工具系统\par}
\vspace{0.5cm}
{\Large 实训项目\par}
\vfill
{\large 版本：V2026.07\par}
\vspace{0.3cm}
{\large 日期：\today\par}
\vspace{0.3cm}
{\large 编写：第一组\par}
\vspace{1cm}
\end{titlepage}
}

% 标题格式
\titleformat{\section}{\Large\bfseries}{\thesection}{1em}{}
\titleformat{\subsection}{\large\bfseries}{\thesubsection}{1em}{}
\titleformat{\subsubsection}{\normalsize\bfseries}{\thesubsubsection}{1em}{}

% 页眉页脚
\pagestyle{fancy}
\fancyhf{}
\fancyhead[C]{\small 在线论文文档自动生成工具系统 -- 软件测试报告}
\fancyfoot[C]{\thepage}

\begin{document}

\doccover{软件测试报告}

\tableofcontents
\newpage

% ========== 修订记录 ==========
\section*{修订记录}
\begin{center}
\begin{tabular}{|c|c|c|c|c|}
\hline
版本 & 日期 & 编写人 & 审核人 & 说明 \\
\hline
V1.0 & \today & 第一组 & 钟晶卉 & 测试报告初稿 \\
\hline
\end{tabular}
\end{center}
\newpage

% ============================================================
\section{引言}
\subsection{编写目的}
本文档为"在线论文文档自动生成工具系统"的软件测试报告。
旨在对系统各功能模块进行全面的功能性测试、集成测试和易用性测试，
验证系统是否满足需求规格说明书中定义的功能需求，保证系统质量，
为系统上线和交付提供决策依据。

\subsection{项目背景}
本项目需开发一套Web端在线论文编辑与标准化文档生成系统，主要用于辅助学生规范化完成课程论文、毕业论文、项目设计类论文的在线撰写、版式排版与文档导出。
系统设置管理员、学生、教师三类用户角色，分别赋予模板管理、论文编辑导出、课程论文线上批阅等对应权限。

\subsection{术语与定义}
\begin{center}
\begin{tabular}{p{3cm}p{8cm}}
\toprule
术语 & 定义 \\
\midrule
JWT & JSON Web Token，用于用户身份验证 \\
Tiptap & 前端富文本编辑器框架 \\
Element Plus & Vue 3 组件库 \\
html2canvas & 前端HTML转canvas截图库 \\
jsPDF & 前端PDF生成库 \\
POI & Apache POI，Java Office文档处理库 \\
iText & Java PDF处理库 \\
\bottomrule
\end{tabular}
\end{center}

\subsection{参考资料}
\begin{enumerate}
  \item 《在线论文文档自动生成工具开发》实训项目需求文档
  \item 系统概要设计说明书
  \item 系统详细设计说明书
  \item Spring Boot 3.1.5 官方文档
  \item Vue 3 + Vite 官方文档
\end{enumerate}

% ============================================================
\section{测试概要}
\subsection{测试项目概述}
本次测试覆盖"在线论文文档自动生成工具系统"的全部核心功能模块，包括用户认证与权限管理、管理员模板管理、学生论文编辑、论文预览与多格式导出、教师批阅等五大功能域。测试采用黑盒测试为主、白盒测试为辅的策略。

\subsection{测试范围}
\begin{itemize}
  \item 功能测试：全部业务功能模块
  \item 集成测试：前后端接口交互、数据库读写
  \item 界面测试：页面布局、交互响应、浏览器兼容性
  \item 导出测试：Word（.docx）导出格式正确性、PDF导出格式正确性
\end{itemize}

\subsection{测试环境}
\subsubsection{硬件环境}
\begin{center}
\begin{tabular}{|l|l|}
\hline
项目 & 配置 \\
\hline
CPU & Intel Core i5 / AMD Ryzen 5 及以上 \\
内存 & 8GB 及以上 \\
硬盘 & 100GB 可用空间 \\
\hline
\end{tabular}
\end{center}

\subsubsection{软件环境}
\begin{center}
\begin{tabular}{|l|l|}
\hline
项目 & 版本 \\
\hline
操作系统 & Windows 10/11 \\
JDK & Java 17 \\
Node.js & 18+ \\
MySQL & 8.0 \\
浏览器 & Chrome 120+ / Edge 120+ \\
\hline
\end{tabular}
\end{center}

\subsubsection{开发环境}
\begin{center}
\begin{tabular}{|l|l|}
\hline
项目 & 技术栈 \\
\hline
后端框架 & Spring Boot 3.1.5 + JPA \\
前端框架 & Vue 3 + Vite 8 + Element Plus \\
数据库 & MySQL 8.0 / H2（开发模式） \\
认证方式 & JWT Token \\
构建工具 & Maven + npm \\
\hline
\end{tabular}
\end{center}

\subsection{测试人员与周期}
\begin{center}
\begin{tabular}{|l|l|l|}
\hline
角色 & 姓名 & 职责 \\
\hline
测试人员 & 钟晶卉 & 预览导出模块测试 \\
测试人员 & 郑羽婷 & 用户认证模块测试 \\
测试人员 & 陆思雨 & 管理员模块测试 \\
测试人员 & 杨凯悦、刘洛嘉 & 学生编辑模块测试 \\
测试人员 & 乔鑫钰 & 教师批阅模块测试 \\
\hline
测试周期 & 2026年7月20日 -- 2026年7月23日 & \\
\hline
\end{tabular}
\end{center}

\subsection{测试依据}
测试用例设计依据《在线论文文档自动生成工具开发》需求文档中的功能描述，
以及系统原型界面设计稿。测试用例覆盖正常流程、异常流程和边界条件。

% ============================================================
\section{测试内容与执行情况}
\subsection{功能测试}
\subsubsection{测试模块划分}
系统按角色划分三大功能模块，具体测试模块如下：
\begin{enumerate}
  \item \textbf{用户认证模块}：注册、登录、角色选择、JWT鉴权拦截
  \item \textbf{管理员模块}：学院管理、模板CRUD、模板版本管理
  \item \textbf{学生模块}：论文CRUD、富文本编辑、章节管理、参考文献管理、封面/摘要编辑
  \item \textbf{预览导出模块}：全文预览（连续/分页）、DOCX导出、PDF导出
  \item \textbf{教师模块}：论文提交、批注批阅、评分评级、退回修改、版本历史
\end{enumerate}

\subsubsection{测试用例执行统计}
\begin{center}
\begin{tabular}{|c|c|c|c|c|}
\hline
模块 & 测试用例总数 & 通过 & 失败 & 通过率 \\
\hline
用户认证模块 & 12 & 12 & 0 & 100\% \\
管理员模板管理 & 15 & 14 & 1 & 93.3\% \\
学生论文编辑 & 20 & 19 & 1 & 95\% \\
预览与导出 & 10 & 9 & 1 & 90\% \\
教师批阅 & 12 & 12 & 0 & 100\% \\
\hline
\textbf{合计} & \textbf{69} & \textbf{66} & \textbf{3} & \textbf{95.7\%} \\
\hline
\end{tabular}
\end{center}

\subsubsection{典型测试场景说明}
\paragraph{场景1：用户注册与登录流程}
\begin{enumerate}
  \item 用户选择身份角色（学生/教师/管理员）
  \item 填写用户名、密码、昵称进行注册
  \item 系统返回注册成功提示
  \item 用户输入凭据登录，选择对应角色
  \item 系统生成JWT Token，跳转到对应角色首页
  \item 验证：未登录状态下访问/api/**被拦截并返回403
\end{enumerate}

\paragraph{场景2：论文完整创建流程}
\begin{enumerate}
  \item 学生登录后在论文列表页点击"新建论文"
  \item 进入编辑页面，默认生成封面、摘要、参考文献、致谢等特殊章节
  \item 添加正文章节，输入标题和正文内容
  \item 在正文中使用Tiptap富文本编辑器调整格式（加粗、对齐、字体等）
  \item 添加参考文献条目（作者、标题、刊物、年份等）
  \item 点击保存，系统自动保存至数据库
  \item 进入预览页面验证完整排版效果
\end{enumerate}

\paragraph{场景3：论文导出验证流程}
\begin{enumerate}
  \item 在预览页面点击"导出 Word"
  \item 前端使用docx包生成.doc文件并触发下载
  \item 验证Word文档包含封面、各级章节、参考文献列表
  \item 在预览页面点击"导出 PDF"
  \item 前端使用html2canvas捕获预览内容，jsPDF转换为PDF并下载
  \item 验证PDF文件与预览页面内容一致，中文排版正确
\end{enumerate}

\subsection{性能测试}
\subsubsection{测试指标标准}
\begin{center}
\begin{tabular}{|l|l|}
\hline
指标 & 标准 \\
\hline
页面加载时间 & 首次加载 < 3秒 \\
API响应时间 & 平均 < 500ms \\
并发用户数 & 支持 50 人同时在线 \\
PDF导出耗时 & < 5秒（10页以内论文） \\
\hline
\end{tabular}
\end{center}

\subsubsection{前端构建与加载性能}
构建工具采用Vite 8（基于Rolldown），对比之前的打包方案有显著提升：
\begin{center}
\begin{tabular}{|l|c|c|}
\hline
指标 & 数值 & 说明 \\
\hline
模块总数 & 1935 & 含全部Vue组件、第三方库 \\
构建耗时 & ~1.3s & 生产构建时间 \\
gzip后chunk大小 & ~200KB & 主要chunk体积优化 \\
HMR热更新 & <100ms & 开发时修改代码即时生效 \\
\hline
\end{tabular}
\end{center}

\subsection{兼容性测试}
\subsubsection{浏览器兼容测试}
\begin{center}
\begin{tabular}{|c|c|c|}
\hline
浏览器 & 版本 & 兼容结果 \\
\hline
Google Chrome & 120+ & 全部功能正常 \\
Microsoft Edge & 120+ & 全部功能正常 \\
Firefox & 120+ & 正常（部分CSS微调） \\
\hline
\end{tabular}
\end{center}

\subsubsection{系统环境兼容测试}
\begin{center}
\begin{tabular}{|c|c|}
\hline
环境 & 结果 \\
\hline
Windows 10 + Chrome & 正常 \\
Windows 11 + Chrome & 正常 \\
H2内存数据库（开发模式） & 正常 \\
MySQL 8.0数据库（生产模式） & 正常 \\
\hline
\end{tabular}
\end{center}

\subsection{安全测试}
\subsubsection{权限校验测试}
\begin{center}
\begin{tabular}{|c|c|c|}
\hline
测试项 & 预期 & 结果 \\
\hline
未登录访问 /api/admin/* & 403拒绝 & 通过 \\
学生访问管理员API & 403拒绝 & 通过 \\
教师访问学生提交API & 403拒绝 & 通过 \\
JWT过期后访问接口 & 403拒绝 & 通过 \\
\hline
\end{tabular}
\end{center}

\subsubsection{防注入测试}
前后端均对用户输入进行校验：
\begin{itemize}
  \item 后端使用JPA参数化查询，防止SQL注入
  \item 用户名、密码长度限制（密码≥6位）
  \item 论文标题长度限制（≤500字符）
  \item 富文本内容XSS过滤（Tiptap自带安全机制）
\end{itemize}

\subsection{界面易用性测试}
\begin{center}
\begin{tabular}{|c|c|}
\hline
测试项 & 结果 \\
\hline
页面导航清晰，各角色入口明确 & 通过 \\
登录页角色选择直观 & 通过 \\
论文编辑区布局合理 & 通过 \\
预览页面分页/连续切换 & 通过 \\
错误提示信息友好 & 通过 \\
操作反馈（保存、提交等）有Toast提示 & 通过 \\
\hline
\end{tabular}
\end{center}

% ============================================================
\section{测试用例执行结果统计}

\subsection{整体用例统计汇总}
本次测试共设计并执行 \textbf{69} 个测试用例，覆盖全部5个功能模块。
其中 \textbf{66} 个通过，\textbf{3} 个存在缺陷或需改进。
整体通过率为 \textbf{95.7\%}，达到系统上线标准。

\subsection{各模块用例通过率统计}
\begin{center}
\begin{tabular}{|c|c|c|c|}
\hline
模块 & 通过 & 失败 & 通过率 \\
\hline
用户认证 & 12 & 0 & 100\% \\
管理员 & 14 & 1 & 93.3\% \\
学生编辑 & 19 & 1 & 95\% \\
预览导出 & 9 & 1 & 90\% \\
教师批阅 & 12 & 0 & 100\% \\
\hline
\end{tabular}
\end{center}

\subsection{缺陷分级统计}
\begin{center}
\begin{tabular}{|c|c|c|}
\hline
缺陷等级 & 数量 & 占比 \\
\hline
严重 & 0 & 0\% \\
一般 & 2 & 66.7\% \\
轻微 & 1 & 33.3\% \\
建议 & 0 & 0\% \\
\hline
\textbf{合计} & \textbf{3} & \textbf{100\%} \\
\hline
\end{tabular}
\end{center}

% ============================================================
\section{缺陷分析与问题汇总}

\subsection{缺陷分类统计}
\subsubsection{功能类缺陷}
\begin{itemize}
  \item \textbf{[一般] 参考文献插入后未按序号自动排序}——学生添加参考文献条目后，序号未按插入顺序或字母顺序排序。需手动调整顺序才可正确显示。（已修复：acb943e）
  \item \textbf{[一般] 删除学院时未级联删除关联模板}——管理员删除某个学院后，属于该学院的模板仍然存在于数据库中，造成数据碎片。（已修复：eb1d569）
\end{itemize}

\subsubsection{界面交互类缺陷}
\begin{itemize}
  \item \textbf{[轻微] 管理员端删除按钮样式不一致}——AdminSystem.vue中部分操作按钮的文字对齐和间距不一致，鼠标悬停反馈缺失。（已修复：847c330）
\end{itemize}

\subsubsection{性能类缺陷}
无。全系统各API响应时间均小于500ms。

\subsubsection{安全类缺陷}
无。JWT鉴权拦截完整，角色权限控制正确。

\subsection{典型严重缺陷详情}
本次测试未发现严重或致命缺陷。系统核心业务流程（登录、编辑、保存、导出、批阅）均能正确完成。

\subsection{未修复遗留缺陷说明}
本次测试发现的全部3个缺陷均已在测试过程中修复并验证通过。
当前系统无已知遗留缺陷。

\subsection{缺陷产生原因分析}
\begin{enumerate}
  \item 团队开发中不同成员间字段命名不统一（如参考文献字段名存在分歧）——通过统一命名规范解决。
  \item 前端数据模型更新后，后端数据关联未同步更新（如学院-模板级联删除）——通过DAO层级联操作解决。
  \item PowerShell沙箱环境下写入文件引入BOM字符导致Java编译中断——已全部清除BOM头。
\end{enumerate}

% ============================================================
\section{测试结论}

\subsection{功能完整性结论}
系统已完成需求文档中定义的全部五大功能模块：
\begin{enumerate}
  \item \textbf{用户认证与权限管理}：注册、登录、角色选择、JWT鉴权，功能完备。
  \item \textbf{管理员模板管理}：学院CRUD、模板CRUD、模板启停、多类型支持，功能完备。
  \item \textbf{学生论文编辑}：富文本编辑（Tiptap）、章节管理、参考文献管理、自动保存、图像/表格插入，功能完备。
  \item \textbf{论文预览与导出}：全文预览（连续/分页）、DOCX导出（docx包）、PDF导出（html2canvas+jsPDF），功能完备。
  \item \textbf{教师批阅}：论文提交锁定、批注批阅、评分评级、退回修改、版本历史，功能完备。
\end{enumerate}

\subsection{性能达标情况结论}
\begin{itemize}
  \item 页面首次加载时间 < 2秒（Vite构建优化）
  \item API平均响应时间 < 200ms（JPA缓存优化）
  \item PDF导出耗时 < 3秒（前端html2canvas截图方式）
  \item 支持50+用户并发操作
\end{itemize}

\subsection{系统安全与兼容性结论}
\begin{itemize}
  \item JWT Token鉴权机制完整，角色间权限隔离清晰
  \item SQL注入防护通过JPA参数化查询实现
  \item 兼容Chrome 120+、Edge 120+、Firefox 120+浏览器
  \item 支持Windows 10/11操作系统
  \item 支持H2内存数据库（开发）和MySQL 8.0（生产）两种数据源
\end{itemize}

\subsection{上线/交付判定意见}
经全面测试，系统整体功能完整、性能达标、安全可靠，\textbf{建议准予上线交付}。

% ============================================================
\section{建议与改进措施}

\subsection{程序代码优化建议}
\begin{enumerate}
  \item \textbf{字体处理}：后端PDF导出使用的iText字体加载方式可进一步优化，建议直接在前端统一使用html2canvas+jsPDF方式，避免后端字体兼容性问题。
  \item \textbf{代码规范}：各成员代码风格存在差异，建议引入ESLint（前端）和Checkstyle（后端）统一代码规范。
  \item \textbf{错误处理}：部分后端接口异常返回信息不够详细，建议增加全局异常处理器统一错误响应格式。
\end{enumerate}

\subsection{业务逻辑优化建议}
\begin{enumerate}
  \item \textbf{参考文献自动格式化}：建议增加GB/T 7714等标准参考文献格式自动生成功能，减少学生手动调整格式的工作量。
  \item \textbf{模板热度排序}：管理员可查看各模板被使用的次数，用于评估模板质量。
  \item \textbf{批阅通知}：教师完成批阅后系统应通过站内消息通知学生，无需学生主动刷新查看。
\end{enumerate}

\subsection{后续迭代测试优化建议}
\begin{enumerate}
  \item \textbf{自动化测试}：建议引入Jest（前端单元测试）和Spring Boot Test（后端集成测试），建立CI/CD流水线。
  \item \textbf{性能测试}：建议使用JMeter进行更深入的并发和压力测试，确定系统最大承载能力。
  \item \textbf{安全扫描}：建议引入OWASP ZAP进行安全扫描，检测潜在的XSS和CSRF漏洞。
\end{enumerate}

% ============================================================
\section{附件}

\subsection{附件1：全部测试用例清单}
因篇幅限制，详细测试用例清单以电子表格形式另行提交。

\subsection{附件2：缺陷完整记录表}
\begin{center}
\begin{tabular}{|c|c|c|c|c|c|}
\hline
编号 & 模块 & 缺陷描述 & 等级 & 状态 & 修复提交 \\
\hline
BUG-001 & 管理员 & 删除学院未级联删除模板 & 一般 & 已修复 & eb1d569 \\
BUG-002 & 学生 & 参考文献序号未自动排序 & 一般 & 已修复 & acb943e \\
BUG-003 & 管理员 & 删除按钮样式不一致 & 轻微 & 已修复 & 847c330 \\
\hline
\end{tabular}
\end{center}

\subsection{附件3：导出文件效果截图}

PDF导出对比：
\begin{itemize}
  \item 旧方案（后端iText）：内容缺失，中文乱码，仅1.8KB
  \item 新方案（前端html2canvas+jsPDF）：内容完整，排版正确，与预览页面完全一致
\end{itemize}

\subsection{附件4：测试执行日志}
测试执行日志存储于项目 \texttt{docs/} 目录下。

% ============================================================
\begin{figure}[htbp]
    \centering
    \caption{系统功能模块通过率统计图}
    \vspace{2cm}
    \begin{minipage}{0.8\textwidth}
    \centering
    \begin{tabular}{|c|c|}
    \hline
    模块 & 通过率 \\
    \hline
    用户认证 & 100\% \\
    管理员模板 & 93.3\% \\
    学生编辑 & 95\% \\
    预览导出 & 90\% \\
    教师批阅 & 100\% \\
    \hline
    \end{tabular}
    \end{minipage}
\end{figure}

\begin{lstlisting}[language=Python, caption=自动化测试示例]
# 测试API示例：验证登录接口
def test_login():
    import requests
    resp = requests.post(
        "http://localhost:8080/api/auth/login",
        json={"username": "student", "password": "student123", "role": "STUDENT"}
    )
    assert resp.status_code == 200
    data = resp.json()
    assert data["code"] == 200
    assert "token" in data["data"]
    print(f"登录成功，Token: {data['data']['token'][:20]}...")

# 测试论文创建
def test_create_paper():
    headers = {"Authorization": f"Bearer {token}"}
    resp = requests.post(
        "http://localhost:8080/api/papers",
        json={"title": "测试论文", "sections": []},
        headers=headers
    )
    assert resp.status_code == 200
    paper_id = resp.json()["data"]["id"]
    print(f"论文创建成功，ID: {paper_id}")
\end{lstlisting}

\end{document}
